% AUTO-GENERATED FILE. Source: pub/claims.toml
Los modelos de lenguaje (LLM) se usan como jueces de la calidad de las
explicaciones, lo que supone que jueces distintos las punt\'uan igual.
Tres jueces LLM puntuaron 192 explicaciones (SHAP, LIME, Anchors, DiCE; tres conjuntos de datos tabulares) en siete
dimensiones con una r\'ubrica de 1 a 5, en dos paneles. El segundo vari\'o el
\emph{prompt}, repiti\'o cada llamada tres veces, puntu\'o los casos de nuevo sin
m\'etricas ni resultado, y puntu\'o 79 explicaciones
SHAP y LIME empeoradas de dos maneras conocidas. Con todos los casos, la correlaci\'on
intraclase ICC(1,1) fue como m\'aximo moderada (m\'aximos:
0.601 y 0.731), y hasta el
67\% de la varianza estaba entre m\'etodos. En las
explicaciones SHAP y LIME con alg\'un peso no nulo el m\'aximo fue
0.41; en SHAP ninguna dimensi\'on
super\'o 0.22, porque los casos recibieron casi la misma
puntuaci\'on. Al quitar siete de diez elementos, todos los jueces bajaron la completitud
y su ICC subi\'o de 0.18 a
0.60. Al asignar los pesos
a variables equivocadas, uno de tres jueces baj\'o la plausibilidad sem\'antica y la
concordancia no subi\'o. Con las m\'etricas en el \emph{prompt}, un juez nombr\'o una
m\'etrica en el 80\% de sus respuestas; sin
ellas puntu\'o m\'as alto y la concordancia en utilidad para auditor\'ia baj\'o de
0.66 a
0.42. Los jueces de esta clase coincidieron en
el m\'etodo y en la forma visible de la explicaci\'on, no en si era correcta. Un solo juez
no deber\'ia usarse como medida confirmatoria, y conviene informar la concordancia por
m\'etodo.
